\documentclass[11pt]{article}
\usepackage[margin=1in]{geometry}
\usepackage{amsmath,amssymb}
\usepackage[hidelinks]{hyperref}

\begin{document}

\section*{Human Review: Divergent orthoprojections in a strictly convex Banach space}

\noindent Original automatic proof: \texttt{solution\_packet.pdf}

\noindent Checked date: 18 June 2026

\noindent Status: Manually checked.

\noindent Verdict: The proof appears correct.

\section*{Human Comments}

The proposed construction appears to give a valid answer to the question in
Badea-Lyubich asking whether divergent iterates of a product of
orthoprojections can occur in a strictly convex Banach space.

The idea is to start from the elementary two-dimensional obstruction and then
blow it up to infinite dimensions.  For each \(0<\rho<1\), the proof builds
a two-dimensional strictly convex block \(E_\rho\) with two norm-one
projections \(P_\rho,Q_\rho\) such that \(P_\rho Q_\rho\) has eigenvalue
\(-\rho^2\).  Taking \(\rho_n\uparrow 1\) gives a sequence of blocks whose
products behave more and more like multiplication by \(-1\).

\section*{Strict Convexity of the Sum}

The key point in the infinite-dimensional construction is the renorming.  One
starts with the \(\ell_\infty\)-sum
\[
  Y=\left(\bigoplus_{n=1}^\infty E_n\right)_{\ell_\infty},
\]
and defines
\[
  \|x\|_X^2
  =
  \|x\|_\infty^2
  +
  \sum_{n=1}^\infty 2^{-n}\|x_n\|_{E_n}^2.
\]
This norm is equivalent to the original \(\ell_\infty\)-sum norm, since
\[
  \|x\|_\infty
  \leq
  \|x\|_X
  \leq
  \sqrt{2}\,\|x\|_\infty.
\]

The strict convexity comes from the weighted \(\ell_2\)-term.  Since each
block \(E_n\) is strictly convex, the function
\[
  x_n\mapsto \|x_n\|_{E_n}^2
\]
is strictly convex on \(E_n\).  Hence
\[
  x\mapsto
  \sum_{n=1}^\infty 2^{-n}\|x_n\|_{E_n}^2
\]
is strictly convex on \(Y\): if \(x\neq y\), then \(x_k\neq y_k\) for some
coordinate \(k\), and the positive weight \(2^{-k}\) preserves the strict
inequality coming from that block.  The term \(x\mapsto \|x\|_\infty^2\) is
convex, so their sum \(\|x\|_X^2\) is strictly convex.  Therefore the norm
\(\|\cdot\|_X\) is strictly convex.

This also explains why the construction is not uniformly convex.  The weights
\(2^{-n}\) tend to zero, so the added curvature becomes arbitrarily weak in
far-out coordinates.  Thus the renorming gives strict convexity without
destroying the asymptotic \(\ell_\infty\)-sum behaviour needed for the
counterexample.

\section*{Contractivity and Divergence}

The blockwise projections
\[
  P(x_n)=(P_nx_n),
  \qquad
  Q(x_n)=(Q_nx_n)
\]
remain contractive for the new norm.  Indeed, they are contractive on each
block, hence contractive for both the \(\ell_\infty\)-term and the weighted
\(\ell_2\)-term.  Therefore they are norm-one projections for
\(\|\cdot\|_X\).

The divergence calculation is then straightforward.  If \(u_n\) is the
distinguished eigenvector in the \(n\)-th block and \(x=(u_n)_n\), then for
\(T=PQ\),
\[
  T^m x
  =
  \left((-\rho_n^2)^m u_n\right)_n.
\]
Since \(\rho_n\uparrow 1\), for every fixed \(m\) the far-out coordinates make
\(\rho_n^{4m}(1+\rho_n^2)\) arbitrarily close to \(2\).  Thus
\[
  \|T^{2m}x-T^{2m+1}x\|_\infty=2,
\]
and hence
\[
  \|T^{2m}x-T^{2m+1}x\|_X\geq 2.
\]
So the orbit is not Cauchy, and the iterates \((PQ)^m\) do not converge
strongly.

\section*{Review Outcome}

The proof appears correct.  The essential mechanism is clear: use
two-dimensional almost-\(\ell_\infty^2\) blocks with product eigenvalues
\(-\rho_n^2\), take \(\rho_n\to1\), and place the blocks in an
\(\ell_\infty\)-sum equipped with a small weighted \(\ell_2\)-correction.  The
correction is enough to force strict convexity, but not enough to remove the
far-out oscillations responsible for divergence.

One minor point to keep in mind is the scalar field.  The proof as written is
most direct over the real field.  If the intended reading of the source
question is complex-linear, the balanced complex variant of the two-dimensional
block should be checked explicitly.  Apart from this standard-field issue, the
argument appears sound.

\end{document}
